\begin{sol}[Péndulo simple vía formulación Lagrangiana]

\subsubsection*{Pregunta 1: Grados de libertad}

La partícula está restringida a moverse sobre una circunferencia de radio $\ell$ (vínculo holonómico $x^2+y^2=\ell^2$), por lo que el sistema tiene $1$ grado de libertad. La coordenada generalizada natural es $q = \theta$, el ángulo respecto a la vertical.

\subsubsection*{Pregunta 2: Lagrangiano}

En coordenadas polares centradas en $O$, con $\theta$ medido desde la vertical:
$$x = \ell\sin\theta, \qquad y = -\ell\cos\theta$$

$$\dot{x} = \ell\dot\theta\cos\theta, \qquad \dot{y} = \ell\dot\theta\sin\theta$$

Energía cinética:
$$T = \frac{1}{2}m(\dot{x}^2+\dot{y}^2) = \frac{1}{2}m\ell^2\dot\theta^2\left(\cos^2\theta+\sin^2\theta\right) = \frac{1}{2}m\ell^2\dot\theta^2$$

Energía potencial gravitatoria, con referencia en $O$:
$$V = mgy = -mg\ell\cos\theta$$

El Lagrangiano es:
$$\boxed{L = T - V = \frac{1}{2}m\ell^2\dot\theta^2 + mg\ell\cos\theta}$$

\subsubsection*{Pregunta 3: Ecuación de Euler-Lagrange}

$$\frac{d}{dt}\left(\frac{\partial L}{\partial\dot\theta}\right) - \frac{\partial L}{\partial\theta} = 0$$

Calculando cada término:
$$\frac{\partial L}{\partial\dot\theta} = m\ell^2\dot\theta \quad\Rightarrow\quad \frac{d}{dt}\left(\frac{\partial L}{\partial\dot\theta}\right) = m\ell^2\ddot\theta$$

$$\frac{\partial L}{\partial\theta} = -mg\ell\sin\theta$$

Sustituyendo:
$$m\ell^2\ddot\theta - (-mg\ell\sin\theta) = 0$$

$$m\ell^2\ddot\theta + mg\ell\sin\theta = 0$$

$$\boxed{\ddot\theta + \frac{g}{\ell}\sin\theta = 0}$$

\paragraph{Comparación con el enfoque newtoniano} Este resultado coincide exactamente con el obtenido aplicando $\vec{F}=m\vec{a}$ en la dirección tangencial: $-mg\sin\theta = m\ell\ddot\theta$. La ventaja del enfoque Lagrangiano es que la tensión de la cuerda —una fuerza de restricción que no realiza trabajo— nunca aparece explícitamente en el cálculo, ya que $\theta$ es la única coordenada independiente compatible con el vínculo.

\end{sol}
